\documentclass[10pt,a4paper]{amsart}
\usepackage{amsmath,amssymb,amsthm,graphicx}
\usepackage{verbatim}
\newenvironment{reference}{\comment}{\endcomment}
%\newenvironment{reference}{}{}
\newenvironment{slogan}{\comment}{\endcomment}
\newenvironment{history}{\comment}{\endcomment}

% For commutative diagrams we use Xy-pic
\usepackage[all]{xy}

% We use 2cell for 2-commutative diagrams.
\xyoption{2cell}
% Theorem environments.
%
\theoremstyle{plain}
\newtheorem{theorem}[subsection]{Theorem}
\newtheorem{proposition}[subsection]{Proposition}
\newtheorem{lemma}[subsection]{Lemma}

\theoremstyle{definition}
\newtheorem{definition}[subsection]{Definition}
\newtheorem{example}[subsection]{Example}
\newtheorem{exercise}[subsection]{Exercise}
\newtheorem{situation}[subsection]{Situation}

\theoremstyle{remark}
\newtheorem{remark}[subsection]{Remark}
\newtheorem{remarks}[subsection]{Remarks}

\numberwithin{equation}{subsection}

% Macros
%
\def\lim{\mathop{\mathrm{lim}}\nolimits}
\def\colim{\mathop{\mathrm{colim}}\nolimits}
\def\Spec{\mathop{\mathrm{Spec}}}
\def\Hom{\mathop{\mathrm{Hom}}\nolimits}
\def\Ext{\mathop{\mathrm{Ext}}\nolimits}
\def\SheafHom{\mathop{\mathcal{H}\!\mathit{om}}\nolimits}
\def\SheafExt{\mathop{\mathcal{E}\!\mathit{xt}}\nolimits}
\def\Sch{\mathit{Sch}}
\def\Mor{\mathop{\mathrm{Mor}}\nolimits}
\def\Ob{\mathop{\mathrm{Ob}}\nolimits}
\def\Sh{\mathop{\mathit{Sh}}\nolimits}
\def\NL{\mathop{N\!L}\nolimits}
\def\CH{\mathop{\mathrm{CH}}\nolimits}
\def\proetale{{pro\text{-}\acute{e}tale}}
\def\etale{{\acute{e}tale}}
\def\QCoh{\mathit{QCoh}}
\def\Ker{\mathop{\mathrm{Ker}}}
\def\Im{\mathop{\mathrm{Im}}}
\def\Coker{\mathop{\mathrm{Coker}}}
\def\Coim{\mathop{\mathrm{Coim}}}

% Boxtimes
%
\DeclareMathSymbol{\boxtimes}{\mathbin}{AMSa}{"02}

%
% Macros for moduli stacks/spaces
%
\def\QCohstack{\mathcal{QC}\!\mathit{oh}}
\def\Cohstack{\mathcal{C}\!\mathit{oh}}
\def\Spacesstack{\mathcal{S}\!\mathit{paces}}
\def\Quotfunctor{\mathrm{Quot}}
\def\Hilbfunctor{\mathrm{Hilb}}
\def\Curvesstack{\mathcal{C}\!\mathit{urves}}
\def\Polarizedstack{\mathcal{P}\!\mathit{olarized}}
\def\Complexesstack{\mathcal{C}\!\mathit{omplexes}}
% \Pic is the operator that assigns to X its picard group, usage \Pic(X)
% \Picardstack_{X/B} denotes the Picard stack of X over B
% \Picardfunctor_{X/B} denotes the Picard functor of X over B
\def\Pic{\mathop{\mathrm{Pic}}\nolimits}
\def\Picardstack{\mathcal{P}\!\mathit{ic}}
\def\Picardfunctor{\mathrm{Pic}}
\def\Deformationcategory{\mathcal{D}\!\mathit{ef}}
\usepackage{cleveref}
\crefname{originalreference}{original reference}{original references}
\Crefname{originalreference}{Original reference}{Original references}
\makeatletter
\newlabel{spaces-lemma-space-presentation}{{0262}{source}{}{sourcepointer-1}{}}
\newlabel{spaces-lemma-space-presentation@cref}{{[originalreference][0][]0262}{[1][0][]source}{}{}{}}
\makeatother
\crefname{originalreference}{original reference}{original references}
\Crefname{originalreference}{Original reference}{Original references}
\makeatletter
\newlabel{morphisms-lemma-etale-over-field}{{02GL}{source}{}{sourcepointer-2}{}}
\newlabel{morphisms-lemma-etale-over-field@cref}{{[originalreference][0][]02GL}{[1][0][]source}{}{}{}}
\makeatother
\crefname{originalreference}{original reference}{original references}
\Crefname{originalreference}{Original reference}{Original references}
\makeatletter
\newlabel{groupoids-proposition-finite-flat-equivalence}{{03BM}{source}{}{sourcepointer-3}{}}
\newlabel{groupoids-proposition-finite-flat-equivalence@cref}{{[originalreference][0][]03BM}{[1][0][]source}{}{}{}}
\makeatother
\title{A quasi-separated space surjected etale by the spectrum of a field is the spectrum of a finite separable extension}
\author{The Stacks Project Authors}
\date{stacks-project:spaces-properties.tex:local-snapshot-20261012; GNU FDL 1.2 or later}
\begin{document}
\maketitle
\noindent\textit{Editorial note.} All delivered mathematical text is the author's own, copied verbatim from the pinned Stacks Project chapter file spaces-properties.tex; no statement or proof step was rewritten, reordered or removed.\par\medskip
\noindent\textit{Editorial note.} Span recomputation by label pairing (LANE\_BRIEF 10.2): in this chapter the author's proof environment follows the statement environment as a standalone block; the spec covers the complete environment pair (statement span, proof span) as recorded in the pool candidate.\par\medskip
\noindent\textit{Editorial note.} Occupation scan: combined claim\_key + statement\_anchor + generic-line scan over all registration files (active and withdrawn) found 13 S10-geom published cards for spaces-properties.tex (prop-occupied.json, 2026-10-14); those 13 labels are excluded from this batch by design. For the 77 free labels: 0 occupied. The registry mentions field (7557 files) and focused-seeds (645) both point to the same S10-geom cards (dual-source check), no other occupation.\par\medskip
\ifdefined\hypertarget\hypertarget{sourcepointer-1}{}\fi
\noindent\textit{Source reference.} Original reference 0262: sources/stacks-chow/tags.txt, source lines 2835--2835. Cross-chapter standard prerequisite cited by the author inside the counted span; it is used as a pointer and its own proof is not part of this item.\par\medskip
\ifdefined\hypertarget\hypertarget{sourcepointer-2}{}\fi
\noindent\textit{Source reference.} Original reference 02GL: sources/stacks-chow/tags.txt, source lines 3232--3232. Cross-chapter standard prerequisite cited by the author inside the counted span; it is used as a pointer and its own proof is not part of this item.\par\medskip
\ifdefined\hypertarget\hypertarget{sourcepointer-3}{}\fi
\noindent\textit{Source reference.} Original reference 03BM: sources/stacks-chow/tags.txt, source lines 4353--4353. Cross-chapter standard prerequisite cited by the author inside the counted span; it is used as a pointer and its own proof is not part of this item.\par\medskip
\begin{lemma}
\label{lemma-point-like-spaces}
Let $S$ be a scheme. Let $k$ be a field.
Let $X$ be an algebraic space over $S$ and assume that there exists
a surjective \'etale morphism $\Spec(k) \to X$.
If $X$ is quasi-separated, then $X \cong \Spec(k')$
where $k/k'$ is a finite separable extension.
\end{lemma}

\begin{proof}
Set $R = \Spec(k) \times_X \Spec(k)$, so that we have a
fibre product diagram
$$
\xymatrix{
R \ar[r]_-s \ar[d]_-t & \Spec(k) \ar[d] \\
\Spec(k) \ar[r] & X
}
$$
By
Spaces, Lemma \ref{spaces-lemma-space-presentation}
we know $X = \Spec(k)/R$ is the quotient sheaf.
Because $\Spec(k) \to X$ is \'etale, the morphisms $s$ and $t$ are
\'etale. Hence $R = \coprod_{i \in I} \Spec(k_i)$ is a disjoint
union of spectra of fields, and both $s$ and $t$
induce finite separable field extensions $s, t : k \subset k_i$, see
Morphisms, Lemma \ref{morphisms-lemma-etale-over-field}.
Because
$$
R = \Spec(k) \times_X \Spec(k)
= (\Spec(k) \times_S \Spec(k)) \times_{X \times_S X, \Delta} X
$$
and since $\Delta$ is quasi-compact by assumption we conclude that
$R \to \Spec(k) \times_S \Spec(k)$ is quasi-compact.
Hence $R$ is quasi-compact as $\Spec(k) \times_S \Spec(k)$ is
affine. We conclude that $I$ is finite. This implies
that $s$ and $t$ are finite locally free morphisms. Hence by
Groupoids, Proposition \ref{groupoids-proposition-finite-flat-equivalence}
we conclude that $\Spec(k)/R$ is
represented by $\Spec(k')$, with $k' \subset k$ finite locally free
where
$$
k' = \{x \in k \mid s_i(x) = t_i(x)\text{ for all }i \in I\}
$$
It is easy to see that $k'$ is a field.
\end{proof}

\end{document}
